%\documentclass[prodmode]{ec-acmsmall-mod}
\documentclass[prodmode]{acmsmall-ec15}

\doi{2764468.2764516}

%\usepackage[small,compact]{titlesec}
\usepackage{amsmath,amssymb,mathtools,dsfont,enumitem,appendix,comment,subfigure}

% Package to generate and customize Algorithm as per ACM style
\usepackage[ruled]{algorithm2e}
\renewcommand{\algorithmcfname}{ALGORITHM}
\SetAlFnt{\small}
\SetAlCapFnt{\small}
\SetAlCapNameFnt{\small}
\SetAlCapHSkip{0pt}
\IncMargin{-\parindent}

% Metadata Information
\acmVolume{X}
\acmNumber{X}
\acmArticle{X}
\acmYear{2015}
\acmMonth{2}
\usepackage[numbers]{natbib}


\DeclareSymbolFont{bbold}{U}{bbold}{m}{n}
\DeclareSymbolFontAlphabet{\mathbbold}{bbold}

\newcommand{\M}{{M_0}}

\newcommand{\eps}{{\epsilon}}
\newcommand{\ind}{{\mathbbold{1}}}
\newcommand{\Adv}{{\mathbf{Adv}}}       
\newcommand{\p}{{\mathbf{P}}}
\newcommand{\pb}{{\mathbf{p}}}           
\newcommand{\vb}{{\mathbf{v}}}           
\newcommand{\F}{{\mathbf{F}}}       
\newcommand{\q}{{\mathbf{Q}}}       
\newcommand{\x}{{\mathbf{x}}}       
\newcommand{\R}{{\mathbb{R}}}       
\newcommand{\opt}{{\mathbf{OPT}}}       
\newcommand{\prp}{{\mathrm{prp}}}    
\newcommand{\calK}{{\cal K}}
\newcommand{\outputs}{{\Rightarrow}}                
\newcommand{\getsr}{{\:\stackrel{{\scriptscriptstyle\hspace{0.2em}\$}}{\leftarrow}\:}}
\newcommand{\andthen}{{\::\;\;}}    % \, \: \; for thinspace, medspace, thickspace
\newcommand{\Rand}[1]{{\mathrm{Rand}[{#1}]}}       % A command with one argument
\newcommand{\Perm}[1]{{\mathrm{Perm}[{#1}]}}       
\newcommand{\Randd}[2]{{\mathrm{Rand}[{#1},{#2}]}} % and with two arguments
\newcommand{\e}{\ensuremath{\mathbf{E}}}

\newcommand{\diag}{\ensuremath{\operatorname{diag}}}
\newcommand{\exit}{\ensuremath{\textsf{exit}}}

%\newtheorem{theorem}{Theorem}
%\newtheorem{prop}[theorem]{Proposition}
%\newtheorem{corollary}[theorem]{Corollary}
%\newtheorem{lemma}[theorem]{Lemma}
%\newtheorem{definition}{Definition}
%\newtheorem{exam}{Example}


\def\EE{{\mathbb{E}}}
\def\PP{{\mathbb{P}}}
\def\BB{{\mathcal{B}}}
\def\MI{{\mathcal{I}}}
\def\setR{\mathbb{R}}
\def\setZ{\mathbb{Z}}
\def\setC{\mathbb{C}}
\def\setN{\mathbb{N}}
\def\bZ{\mathbf{Z}}
\def\bz{\mathbf{z}}
\def\FF{\mathcal{F}}
\newcommand{\fall}{\,\forall\,}
\def\vecV{\underbar{v}}

\usepackage[colorlinks,
            linkcolor=blue,
            citecolor=blue,
            urlcolor=magenta,
            linktocpage,
            plainpages=false,
            hypertexnames=false,
	    	pdfauthor={Siddhartha Banerjee},
%            pagebackref=true,
	    	pdftex]{hyperref}
%\usepackage{nameref}
%\usepackage{bookmark}

\usepackage{color}
\newcommand{\todo}{\color{red}}


\clubpenalty=10000 
\widowpenalty = 10000


%%%%%%%%%%%%%%%%%%%%%%%%%%%%%%%%%%%%%%%%%%%%%%%%%%%%%%%%%%%%%%%%%%%%%%%%%%%
\begin{document}

% Page heads
\markboth{Horton, Johari}{Eliciting Buyer Preferences as a Market Design Problem}

% Title portion
%\title{A Ride Whenever you Need One: The Value of Dynamic Pricing in Ride-sharing Platforms}
\title{At What Quality and What Price? Eliciting Buyer Preferences as a Market Design Problem}
\author{JOHN J. HORTON
\affil{Leonard N. Stern School of Business, New York University}
RAMESH JOHARI
\affil{Stanford University}
}

\begin{abstract}
  \noindent  Buyers and sellers in markets often signal to inform the other side about their preferences.
  Both have a mutual incentive to reveal information with respect to horizontal differentiation, but the case of vertical differentiation is more complex:
  a buyer claiming they place a high value on quality may attract more sellers of the right ``type'' increasing efficiency, but they might also simply pay a higher price.
  Although an efficiency-minded social planner may not care about higher prices, if this fear prevents a buyer from stating his of her true preferences, then desirable sorting caused by information-revelation may be unattainable.
  In this paper, we consider the buyer's vertical differentiation disclosure problem through the lens of a large field experiment conducted in an online labor market. 
  A new signaling mechanism was introduced into the market that allowed buyers to state their relative preferences over price and quality.
  We find that the buyer signal improved seller-side sorting, with more sellers going to buyers of the right ``type''; the total number of applications also fell.
  However, sellers also clearly tailored their wages bid to the type of buyer they faced.
  Despite this markup, buyers chose to honestly disclose their preferences, suggesting they found the sorting effect to dominate the bargaining power effect.
\end{abstract} 

\category{J4}{Social and Behavioral Sciences}{Economics}

% A category with the (minimum) three required fields
%\category{H.4}{Information Systems Applications}{Miscellaneous}
%\category{C.2.2}{Computer-Communication Networks}{Network Protocols}
\terms{Economics; Experimentation} 

\keywords{signaling, matching, market design}

\begin{bottomstuff}
Author contact information, data sets, code, and the current version of this draft are currently or will be available at \href{http://www.john-joseph-horton.com/}{http://www.john-joseph-horton.com/}.

This work was carried out in part while the authors were employed by oDesk Corporation.
\end{bottomstuff}

\maketitle
%%%%%%%%%%%%%%%%%%%%%%%%%%%%%%%%%%%%%%%%%%%%%%%%%%%%%%%%%%%%%%%%%%%%%%

%\section{Extended Abstract}


%\vspace{-0.35cm}
%Acknowledgments
%\begin{acks}
%Part of this work was carried out when SB was a technical consultant at Lyft.
%We would like to thank the data science team at Lyft, particularly Adam Greenhall, Chris Pouliot, and Chris Sholley; the formulation of our model greatly benefited from their inputs, and also from discussions with Lyft drivers.
%We also thank Anilesh Krishnaswamy for several useful technical discussions.

%We gratefully acknowledge support from the National Science Foundation via grants CMMI-1234955 and CNS-1343253, and the DARPA GRAPHS program.
%\end{acks}


%\bibliographystyle{acmsmall}
%\bibliography{RideShare}

% Electronic Appendix
%\elecappendix
%\input{Appendix.tex}

\end{document}

